\documentclass[12pt,a4paper]{article}
\usepackage{amsmath}
\usepackage{amssymb}  % Provides \mathbb and other symbols
\usepackage{graphicx}
\usepackage{subcaption}
\usepackage{enumitem}
\usepackage{listings}
\usepackage{hyperref}
\usepackage{ulem}
\usepackage{xcolor}
\usepackage{amsthm}
\usepackage{tcolorbox}  % For colored boxes
\usepackage[table]{xcolor} % For coloring table cells
\usepackage{array}         % For better table alignment
\usepackage{mdframed}
\usepackage{cancel}
\usepackage[margin=2.5cm]{geometry}
\setlength{\parindent}{0pt}
\setlength{\parskip}{0.8em}

\begin{document}
\begin{titlepage}
\centering
\vspace*{\fill}
{\Huge Algebra} \\[0.5cm]
{\Huge Home Work 5}\\[1.5cm]
{\Large Student: 26-712-224}\\[0.5cm]
{\Large \today}
\vspace*{\fill}
\end{titlepage}

\newpage
\noindent
A) Let $G$ be a group and $H<G$ a subgroup. We assume that $G/H=H\backslash G$, i.e., every left $H$-coset is also a right $H$-coset and vice versa.
\begin{enumerate}
\renewcommand{\labelenumi}{(\arabic{enumi})}
\item Show that for every $g\in G$, there exists $g'\in G$ such that $g\cdot H=H\cdot g'$.
\item Let $g$ and $g'$ be as in (1). Show that $1\in H\cdot(g'g^{-1})$.
\item Use (2) to deduce that $H$ is a normal subgroup of $G$.
\item As an application, show the following: if $H<G$ is a subgroup such that $\operatorname{card}(G/H)=2$, then $H$ is normal in $G$.
\end{enumerate}

\vspace{1cm}
\noindent
Answer:\\
\\
(1) Every element of $G/H$ has the form $g.H$ for some $g \in G$. Also every element in $H \backslash G$ has the form $H.k$ for
some $k \in G$.  We are given that the two collections of cosets are the same.
This means given $g.H$ in the set of left-cosets there will
be a corresponding right-coset (and vice versa) - which we can write as $H.g'$ for some $g' \in G$.\\
\\
(2) Using the notation from (1): 
\begin{align*}
    g.H = H.g' \implies g.H.g^{-1} = H.(g'g^{-1})
\end{align*}
As $H < G \implies 1 \in H$.
Consequently: 
\begin{align*}
    g1g^{-1} = 1 \in H.(g'g^{-1}) \tag{A.1}
\end{align*}
\\
(3) From (A.1) we have: $1 \in H.(g'g^{-1}) \implies g \in H.g'$.  However, any coset is defined by any of its representative
elements and hence: $H.g = H.g'$.  Taking into account with $g.H = H.g'$ we must conclude that:
\begin{align*}
    g.H = H.g \;\;\;\;\;\;\;\;\; \text{(true $\forall g \in G$)}
\end{align*}
As this must hold in all cases (i.e. for all left and right cosets) we must have that $H \triangleleft G$.\\
\\
(4) If $\text{Card}(G/H)=2$.  We know that the number of left (and right cosets) is 2.\\
$H$ itself will be BOTH a left AND a right coset.  The remaining elements from $G \setminus H$ (set difference) must then reside in the remaining left coset and
also in the remaining right coset.\\
\\
So we have two left cosets: $(H, g.H)$ - for some $g \in G$, and two righ cosets: $(H, H.g')$ - for some $g' \in G$.\\
\\
We, therefore, must have that: $g.H = H.g'$.  The conclusion is that $H \triangleleft G$.

\newpage
\noindent
B) Let $S_3$ be the symmetric group on $3$ elements, i.e., the group of bijections from $\{1,2,3\}$ to itself.
\begin{enumerate}
\renewcommand{\labelenumi}{(\arabic{enumi})}
\item What is the cardinality of $S_3$?
\item Show that the center of $S_3$ is trivial, i.e., that $Z(S_3)=\{\operatorname{id}_{\{1,2,3\}}\}$.
\item Let $\sigma\in S_3$ be the element given by $\sigma(1)=2$, $\sigma(2)=3$ and $\sigma(3)=1$. What is the order of $\sigma$?
\item Using Exercise A), show that $\langle\sigma\rangle$ is normal in $S_3$.
\item Show that $S_3/\langle\sigma\rangle$ is isomorphic to $\mathbb{Z}/2\mathbb{Z}$.
\item Using the previous questions, compute the derived series of $S_3$ and show that $S_3$ is solvable.
\end{enumerate}

\vspace{1cm}
\noindent
Answer:\\
\\
(1) $\text{Card}(S_3) = 3! = 6$ - this is the total number of permutations of 3 elements.\\
\\
(2) The centre of $S_3 = \{s \in S_3 \; | \; g s = s g; \; \forall g \in S_3 \}$.\\
\\
Consider the element $(1 2) \in S_3$ - if $s \in Z(S_3)$ then:
\begin{align*}
    s(1 \; 2) = (1 \; 2)s
\end{align*}
As 3 is "not moved" that means that $s(1) \ne 3$ and $s(2) \ne 3$ and so $s(3)=3$.\\
Application of the same logic shows that: $s(1)=1; s(2)=2$ and hence $s=id$.\\
\\
(3) We are given that:
\begin{align*}
\sigma =
\begin{pmatrix}
1 & 2 & 3 \\
2 & 3 & 1
\end{pmatrix}
\end{align*}
Clearly: $\sigma^3=id$ and hence: $\text{ord}(\sigma)=3$.\\
\\
(4) $\langle \sigma \rangle = \{id, \sigma, \sigma^2 \} < S_3$.\\
\\
Now $\langle \sigma \rangle$ has order 3 and $S_3$ has order 6.  From this we can conclude that
$[S_3 : \langle \sigma \rangle]=2$.\\
\\
Part (6) of Q1 means that $\langle \sigma \rangle \triangleleft S_3$.\\
\\
(5) $S_3 / \langle \sigma \rangle \triangleleft S_3$ is a group with 2 elements.  The only groups of order 2 are
all isomorphic to $\mathbb{Z} \ 2 \mathbb{Z}$.

\newpage
\noindent
(6)
 
\newpage
\noindent
C) The goal of this exercise is to show that the notions of left and right actions are equivalent. We fix a group $G$ and a set $X$.
\begin{enumerate}
\renewcommand{\labelenumi}{(\arabic{enumi})}
\item Let $a:G\times X\to X$ be a left action of $G$ on $X$. Show that the map $b:X\times G\to X$, given by $b(x,g)=a(g^{-1},x)$ for $x\in X$ and $g\in G$, defines a right action of $G$ on $X$.
\item Conversely, let $b:X\times G\to X$ be a right action. Show that the map $a:G\times X\to X$, given by $a(g,x)=b(x,g^{-1})$ for $x\in X$ and $g\in G$, defines a left action of $G$ on $X$.
\item Show that the constructions in (1) and (2) are inverses of each other and define bijections between the sets of right and left actions of $G$ on $X$.
\end{enumerate}

\vspace{1cm}
\noindent
Answer:\\
\\


\end{document}